\subsection{年度研究计划}
研究期限为2027年1月至2028年12月。申请人负责设计、核查和论文组织，\XJField{student-count}名研究生按成员表分工，重点培养2名参与完整研究过程。

\subsubsection{2027年：完成数据核验、尺度分析与概率预测基线}
第一季度对接1个新疆主研究场站的连续24个月记录，核验授权、关键字段与质量。第二季度按15分钟网格和50\%/15\%/15\%/20\%时间划分建库，完成持续性、分位数梯度提升及标准时序基线，核查日照、天气和异常分层。第三季度开展输入与尺度消融，训练DWT-Former主模型，以3个随机种子重复实验。第四季度完成80\%和90\%区间初步校准，保存历史留出预测与完整误差轨迹，核查天气及爬坡样本量。提交样本索引、质量与尺度增益报告、预测程序和中期报告；未获最低要求的新疆实测时，说明地域验证未完成并调整任务。

\subsubsection{2028年：完成天气变化下的校准比较、爬坡检验与成果整理}
第一季度验证30天滚动校准及14天、60天窗口敏感性，锁定模型、告警与统计规则，保存已冻结的天气分组。第二季度比较每次100条同边际轨迹的三种排序，以30分钟下降超过装机容量20\%为主事件，回放独立测试段。第三季度检验50条、200条轨迹及10\%、30\%爬坡阈值的敏感性，整理天气个例、失败案例及统计区间，形成论文初稿与原型。第四季度投稿论文、发布软件，提交报告、使用说明及结题材料，归档2名重点培养研究生的实验记录与总结；如实报告阴性结果和未完成条件。

\subsection{预期研究结果}
预期形成数据处理及评价流程、概率预测与校准实现、爬坡风险离线原型各一套，交付字段规范、质量报告和复现实验索引。方法报告说明气象与尺度增益、校准效率、时间依赖贡献及其失效条件。

完成并投稿论文1篇，交付算法软件1项（含预测、校准及预警模块），\XJField{student-training-goal}。7名研究生参与，2名重点培养对象以实验记录、程序和总结考核，人选及投入据实落实。软件附版本、测试案例与使用说明，推进著作权登记；录用、登记及培养结果据实列报。成果署名、基金标注与档案移交遵照基金规则及授权，最终考核以批准合同为准。
